% fig_pinn_architecture.tex
% TR-46: PINN architecture diagram
% Input → hidden layers → output + physics constraint residuals

\begin{tikzpicture}[
  every node/.style={font=\scriptsize},
  neuron/.style={circle, draw=blue!50, fill=blue!10, minimum size=0.50cm,
                 inner sep=0pt},
  layer/.style={draw=sambpblue, thick, rounded corners=4pt, fill=blue!5,
                minimum width=1.1cm, minimum height=4.0cm, align=center},
  phys/.style={draw=orange!70!black, thick, rounded corners=4pt,
               fill=orange!5, minimum width=2.8cm, minimum height=0.7cm,
               align=center, font=\tiny},
  arrow/.style={-{Stealth[length=5pt]}, thick},
  redarrow/.style={-{Stealth[length=5pt]}, thick, orange!70!black},
]

%% ── Input layer ───────────────────────────────────────────────────────────
\node[neuron, fill=green!15, draw=green!50!black] (in1) at (0,  1.275) {};
\node[neuron, fill=green!15, draw=green!50!black] (in2) at (0,  0.425) {};
\node[neuron, fill=green!15, draw=green!50!black] (in3) at (0, -0.425) {};
\node[neuron, fill=green!15, draw=green!50!black] (in4) at (0, -1.275) {};
\node[left=6pt, font=\tiny] at (in1.west) {$k_\mathrm{ibr}/0.55$};
\node[left=6pt, font=\tiny] at (in2.west) {$\alpha/1.10$};
\node[left=6pt, font=\tiny] at (in3.west) {$R_\mathrm{arc}/0.025$};
\node[left=6pt, font=\tiny] at (in4.west) {fault type};

%% ── Hidden layers (3 × 64 neurons, shown as stacks) ──────────────────────
\node[layer] (h1) at (2.5,  0.0) {\footnotesize 64\\ReLU};
\node[layer] (h2) at (4.5,  0.0) {\footnotesize 64\\ReLU};
\node[layer] (h3) at (6.5,  0.0) {\footnotesize 64\\ReLU};

%% ── Output layer (4 classes) ─────────────────────────────────────────────
\foreach \i/\lab/\col in {
  1/{P(MISS)}/red!40,
  2/{P(87L)}/blue!40,
  3/{P(Z1)}/green!40,
  4/{P(EXT)}/gray!40} {
  \pgfmathsetmacro{\y}{(2.5 - \i) * 0.85}
  \node[neuron, fill=\col, minimum size=0.55cm]
    (out\i) at (8.8, \y) {};
  \node[right=6pt, font=\tiny] at (out\i.east) {\lab};
}

%% ── Softmax ───────────────────────────────────────────────────────────────
\node[draw=gray!50, rounded corners=3pt, fill=gray!5,
      font=\tiny, inner sep=3pt, align=center]
  at (8.8, -2.20) {Softmax\\$\sum P_i=1$};

%% ── Connection arrows ─────────────────────────────────────────────────────
\foreach \i in {1,2,3,4} {
  \draw[arrow, blue!30] (in\i.east) -- (h1.west |- in\i);
}
\draw[arrow, blue!40] (h1.east) -- (h2.west);
\draw[arrow, blue!40] (h2.east) -- (h3.west);
\foreach \i in {1,2,3,4} {
  \draw[arrow, blue!30] (h3.east |- out\i) -- (out\i.west);
}

%% ── Physics constraint boxes ──────────────────────────────────────────────
\node[phys] (c1) at (3.5, -2.80)
  {C1: $k<0.06 \Rightarrow P(\mathrm{MISS})\geq0.90$};
\node[phys] (c2) at (6.0, -2.80)
  {C2: $k\geq0.10,\,\alpha\leq0.80 \Rightarrow P(Z1)+P(87L)\geq0.95$};
\node[phys] (c3) at (9.5, -2.80)
  {C3: $\alpha>1.0 \Rightarrow P(\mathrm{EXT})\geq0.95$};

%% ── Loss function ────────────────────────────────────────────────────────
\node[draw=red!40, fill=red!5, rounded corners=3pt,
      font=\tiny, inner sep=4pt, align=center]
  at (5.5, -4.10)
  {$\mathcal{L} = \underbrace{\mathcal{L}_\mathrm{CE}}_\text{data loss}
   + \lambda_\mathrm{phy}\underbrace{(R_{C1}+R_{C2}+R_{C3})}_\text{physics residual}$
   \quad $\lambda_\mathrm{phy}=0.20$};

%% Arrows from outputs to constraint boxes
\draw[redarrow, dashed] (out1.south) |- (c1.east);
\draw[redarrow, dashed] (out3.south) |- (c2.east);
\draw[redarrow, dashed] (out4.south) |- (c3.west);

%% Arrows from constraints to loss
\draw[redarrow] (c1.south) |- (5.5, -3.80);
\draw[redarrow] (c2.south) -- (5.5, -3.80);
\draw[redarrow] (c3.south) |- (5.5, -3.80);
\draw[redarrow, very thick] (5.5, -3.80) -- (5.5, -3.75);

%% Annotations
\node[font=\small\bfseries, blue!60!black] at (2.5, 2.5) {Layer 1};
\node[font=\small\bfseries, blue!60!black] at (4.5, 2.5) {Layer 2};
\node[font=\small\bfseries, blue!60!black] at (6.5, 2.5) {Layer 3};
\node[font=\small\bfseries, green!50!black] at (0, 2.5) {Input};
\node[font=\small\bfseries, gray!60]        at (8.8, 2.5) {Output};

\end{tikzpicture}
